% BOOK-FILE-ID: f-73b81d4e9a02
\subsubsection{虚数解型を母関数で解く}\label{節：虚数解型を母関数で解く}
本編\sectionref{節：特性方程式の解が虚数の場合}の例を,母関数でも解いてみよう.
\[
  a_1=0,\quad a_2=1,\qquad a_{n+2}=2a_{n+1}-2a_n
\]
に対し,本編の項の番号を保って
\[
  F(x)=\sum_{n=1}^\infty a_nx^{n-1}
\]
と置く.二次以上の係数が漸化式によって消えるため,
\[
  (1-2x+2x^2)F(x)=x
\]
となる.分母を複素数の範囲で因数分解し,部分分数分解すると,
\[
  \begin{aligned}
    F(x)&=\frac{x}{(1-(1+i)x)(1-(1-i)x)}\\
      &=-\frac{i}{2}\frac{1}{1-(1+i)x}
         +\frac{i}{2}\frac{1}{1-(1-i)x}.
  \end{aligned}
\]
等比級数の係数を比べれば,
\[
  a_n=-\frac{i}{2}(1+i)^{n-1}+\frac{i}{2}(1-i)^{n-1}
\]
となり,本編と同じ一般項を得る.
一般にも,$a_{n+2}+pa_{n+1}+qa_n=0$なら,
\[
  F(x)=\frac{a_1+(a_2+pa_1)x}{1+px+qx^2}
\]
である.分母の二つの因子から二つの累乗が現れ,分子に記録された初期値からその係数が決まる.
一般項の形を先に仮定せずに導けたのである.

\subsubsection{付加項がある場合：各回の作用を足す}
次に,$a_{n+2}+pa_{n+1}+qa_n=f_n$を考えよう.
ここでは$p=-2\rho\cos\theta$, $q=\rho^2$, $0<\theta<\pi$とし,$a_1=a_2=0$の場合から始める.
$H(x)=\sum_{n=1}^\infty f_nx^{n-1}$とすれば,
\[
  F(x)=\frac{x^2H(x)}{1+px+qx^2}
\]
である.分母の逆数の係数は,
\[
  \begin{aligned}
    \frac{1}{1+px+qx^2}
      &=\sum_{j=0}^\infty g_jx^j,\\
    g_j&=\rho^j\frac{\sin((j+1)\theta)}{\sin\theta}
  \end{aligned}
\]
となる.実際,$g_0=1$, $g_1=2\rho\cos\theta=-p$であり,それ以降の係数は付加項のない漸化式を満たす.
積の係数を比べると,
\[
  \begin{aligned}
    a_n&=\sum_{k=1}^{n-2}f_kg_{n-k-2}\\
       &=\sum_{k=1}^{n-2}f_k\rho^{n-k-2}
         \frac{\sin((n-k-1)\theta)}{\sin\theta}
         \quad(n\ge3)
  \end{aligned}
\]
を得る.$k$回目に加わった$f_k$がその後に作る振動を,すべて足した式である.
初期値が0でない場合は,同じ初期値を持つ付加項のない解を加えればよい.
このように,母関数の積は外から加わる作用の重ね合わせにも使える.

例えば$\rho=1$, $\theta=\frac{\pi}{2}$で,
\[
  a_1=a_2=0,\qquad
  a_{n+2}+a_n=\cos\frac{(n-1)\pi}{2}
\]
とすると,$0,0,1,0,-2,0,3,0,\ldots$となる.
帰納法で
\[
  a_n=\frac{n-1}{2}\sin\frac{(n-2)\pi}{2}
\]
を確かめられる.付加項のない式では振幅が一定だったが,同じ角の進み方を持つ作用が加わり続けると振幅が増大する.
これは離散的な共振の例である.ここで$n$が現れるのは,付加項のない特性方程式が重解だからではない.
元の特性方程式の解は$\pm i$で異なるが,同じ振動を作る付加項が加わったためである.
